\documentclass[a4paper,11pt]{article}
\usepackage[utf8]{inputenc}
\usepackage{amsmath,amssymb,amsthm,mathtools}
\usepackage{booktabs}
\usepackage{graphicx}
\usepackage{hyperref}
\usepackage{cleveref}
\usepackage[numbers]{natbib}

\newtheorem{theorem}{Teorema}
\newtheorem{lemma}{Lemma}
\newtheorem{definition}{Definizione}

\title{FeatherTex Demo: Articolo Scientifico Multi-File}
\author{Autore Demo}
\date{\today}

\begin{document}
\maketitle

\begin{abstract}
Questo progetto dimostrativo mostra le funzionalità di \textbf{FeatherTex} per ricercatori:
struttura multi-file con \texttt{\textbackslash input}, autocompletamento di \texttt{\textbackslash cite} e
\texttt{\textbackslash cref}, tabelle \texttt{booktabs}, bibliografia BibTeX automatica,
conteggio parole e linter tipografico.
\end{abstract}

% !TeX root = ../main.tex
\section{Introduzione}\label{sec:intro}

La scrittura scientifica in \LaTeX{} beneficia di un flusso di lavoro rapido e modulare~\cite{knuth1984,lamport1994}.
In questo file secondario il commento magico \texttt{\% !TeX root = ../main.tex} permette di
premere \texttt{Ctrl+B} direttamente da qui compilando il documento principale.

\begin{definition}[Spazio Completo]\label{def:completo}
Uno spazio metrico $(X, d)$ si dice completo se ogni successione di Cauchy in $X$ converge ad un elemento di $X$.
\end{definition}

\begin{theorem}[Teorema del Punto Fisso]\label{thm:banach}
Sia $(X, d)$ uno spazio metrico completo non vuoto e sia $T : X \to X$ una contrazione.
Allora $T$ ammette un unico punto fisso $x^* \in X$.
\end{theorem}


% !TeX root = ../main.tex
\section{Risultati e Confronto}\label{sec:results}

Dal~\cref{thm:banach} e dalla~\cref{def:completo} segue la convergenza geometrica dell'iterazione:
\begin{equation}\label{eq:conv}
  d(x_n, x^*) \le \frac{q^n}{1 - q}\, d(x_1, x_0), \qquad q \in (0, 1).
\end{equation}

La~\cref{tab:bench} riassume i tempi tipici di compilazione incrementale e rendering~\cite{vaswani2017}.

\begin{table}[htbp]
  \centering
  \begin{tabular}{lcc}
    \toprule
    \textbf{Fase} & \textbf{Metodo classico} & \textbf{FeatherTex} \\
    \midrule
    Preambolo (.fmt) & 0.90 s & 0.12 s \\
    Rendering PDF (Poppler+Cairo) & 418 ms & 15 ms \\
    Anteprima pre-compilata & --- & Istantanea \\
    \bottomrule
  \end{tabular}
  \caption{Confronto prestazioni nella pipeline di compilazione.}
  \label{tab:bench}
\end{table}



\bibliographystyle{plainnat}
\bibliography{references}

\end{document}

